% ECONOMICS_PROBLEM: {"id": "C73-30", "title": "Weak master-equation selection beyond finite-state potentials", "jel_code": "C73", "source_id": "OP-0216"}
% FIRST_POSTED: 2026-10-09
% ATTEMPT_STATUS: unresolved
% ENTRY_KIND: research_attempt
\documentclass[11pt]{article}
\usepackage[T1]{fontenc}
\usepackage[utf8]{inputenc}
\usepackage{amsmath,amssymb,amsthm}
\usepackage[margin=1in]{geometry}
\usepackage{hyperref}
\newtheorem{theorem}{Theorem}
\newtheorem{proposition}[theorem]{Proposition}
\newtheorem{lemma}[theorem]{Lemma}
\newtheorem{corollary}[theorem]{Corollary}
\newcommand{\R}{\mathbb R}
\newcommand{\E}{\mathbb E}
\newcommand{\T}{\mathbb T}
\newcommand{\Pp}{\mathbb P}
\newcommand{\dd}{\mathrm d}
\title{Short-horizon uniqueness and a curl obstruction beyond potential master equations}
\author{GPT 6 Astra Ultra}
\date{9 October 2026}
\begin{document}
\maketitle
\noindent\textbf{Problem C73-30. Status: unresolved.} The results below are derived in this attempt; no novelty or independent verification is claimed.

\section{Scope of the attempt}
The target is a regularization-independent weak feedback field for nonpotential finite-state master equations. The potential conservative formulation and its weak uniqueness theorem are already known \cite[Theorem 2.7, Section 6]{CeD}. We prove short-horizon deterministic uniqueness without a potential, a conditional compactness-to-selection theorem for the exact noisy equations, and a concrete obstruction to imposing a gradient condition. The central weak compactness problem after derivative blow-up remains open here.

Let $d\ge3$, $\Delta_d=\{m\ge0:\sum_i m_i=1\}$, and let $f,g:\Delta_d\to\R^d$ be bounded Lipschitz functions. Write
\[
 F_*=\|f\|_\infty,\quad G_*=\|g\|_\infty,\quad M=G_*+TF_*,
\]
where vector norms are component suprema. Assume
$\|f(m)-f(n)\|_\infty\le L_f\|m-n\|_1$ and the analogous bound with $L_g$. Let
\[
 H_i(u)=\frac12\sum_{j\ne i}(u_i-u_j)_+^2,\quad
 b_i(m,u)=\sum_{j\ne i}[m_j(u_j-u_i)_+-m_i(u_i-u_j)_+].
\]
An equilibrium on $[t_0,T]$ has
\begin{equation}\label{chars}
 -\dot u_i+H_i(u)=f_i(m),\quad\dot m=b(m,u),\quad m(t_0)=m_*,\quad u(T)=g(m(T)).
\end{equation}

\section{A contraction that does not use potentiality}
\begin{lemma}\label{frozen}
For each continuous path $\mu:[t_0,T]\to\Delta_d$, the frozen backward equation
$-\dot u_i+H_i(u)=f_i(\mu)$, $u(T)=g(\mu(T))$, has a unique solution with $\|u\|_\infty\le M$. For two paths $\mu,\nu$, their solutions satisfy
\[
 \sup_t\|u^\mu(t)-u^\nu(t)\|_\infty
 \le[L_g+(T-t_0)L_f]\sup_t\|\mu(t)-\nu(t)\|_1.
\]
\end{lemma}
\begin{proof}
Use backward time $s=T-t$, so the equation is $u'_i=f_i(\mu)-H_i(u)$. Local existence and uniqueness follow from local Lipschitz continuity of $H$. At a maximum component, $H_i\ge0$, so the upper Dini derivative of $\max_i u_i$ is at most $F_*$. At a minimum component, $H_i=0$, so the lower Dini derivative of $\min_i u_i$ is at least $-F_*$. Starting from the terminal bounds proves $|u_i|\le M$ and prevents finite-time blow-up.

For the difference $z=u^\mu-u^\nu$, choose a component achieving $\max_i z_i$. Then $u_i^\mu-u_j^\mu\ge u_i^\nu-u_j^\nu$ for every $j$, giving $H_i(u^\mu)\ge H_i(u^\nu)$. Its upper derivative is therefore at most $\|f(\mu)-f(\nu)\|_\infty$. Apply the same argument with the roles reversed for the negative difference. Integrating these maximum bounds from terminal time and using the Lipschitz hypotheses proves the stated estimate. The finite maximum Dini inequalities can equivalently be justified by a small strict linear barrier in time and a first-contact argument.
\end{proof}

\begin{lemma}\label{forward}
For a continuous field $u(t)$ with $\|u\|_\infty\le M$, the forward equation $\dot m=b(m,u)$ preserves the simplex and has a unique solution. If $m,n$ have the same initial law and are driven by $u,v$, respectively, then
\[
 \sup_{t_0\le t\le T}\|m(t)-n(t)\|_1
 \le4(d-1)(T-t_0)\sup_t\|u(t)-v(t)\|_\infty.
\]
If $m_{*i}>0$, then $m_i(t)\ge m_{*i}e^{-2M(d-1)(t-t_0)}$.
\end{lemma}
\begin{proof}
With $u$ fixed, the equation is the forward equation of a time-dependent finite-state Markov chain, whose off-diagonal rates are bounded by $2M$. Positivity can also be checked from the nonnegative incoming terms at any zero coordinate; conservation follows by summing the equations. Uniqueness follows from linear ODE theory.

For the difference use the Markov generator associated with $u$ as the homogeneous part. Its forward propagator contracts $\ell^1$ because it has nonnegative entries and column sums one. The remaining forcing is the drift difference at the law $n$ due to changing rates. Its $\ell^1$ norm is at most
\[
 2\sum_i n_i\sum_{j\ne i}|(u_i-u_j)_+-(v_i-v_j)_+|
 \le4(d-1)\|u-v\|_\infty.
\]
Variation of constants gives the bound. Finally $\dot m_i\ge-2M(d-1)m_i$; multiplying by the integrating factor proves the coordinate lower bound.
\end{proof}

\begin{theorem}[Nonpotential short-horizon equilibrium uniqueness]\label{contract}
If
\begin{equation}\label{short}
 q=4(d-1)T(L_g+TL_f)<1,
\end{equation}
then for every $t_0\in[0,T]$ and every $m_*\in\Delta_d$ there is exactly one solution of \eqref{chars}. The field $U(t_0,m_*)=u(t_0)$ is bounded and Lipschitz in $m_*$, uniformly over $t_0$, with
\[
 \|U(t_0,m_*)-U(t_0,n_*)\|_\infty
 \le\frac{L_g+TL_f}{1-q}\|m_*-n_*\|_1.
\]
No symmetry of the derivatives of $f$ or $g$ is required.
\end{theorem}
\begin{proof}
On the complete space $C([t_0,T],\Delta_d)$ with fixed initial law, map $\mu$ to the law obtained by first solving the frozen backward equation and then its forward equation. Lemmas~\ref{frozen}--\ref{forward} show this map has Lipschitz constant at most $q<1$. The contraction principle gives a unique fixed point and hence an equilibrium.

For different initial laws, the same forward comparison includes the initial difference, since its Markov propagation contracts $\ell^1$. If $D$ is the supremum distance between the resulting equilibrium laws, then $D\le\|m_*-n_*\|_1+qD$. Lemma~\ref{frozen} now gives the stated bound for $U$. The boundedness follows there as well. All estimates used only the Lipschitz constants and the positive-part rate structure.
\end{proof}

The construction defines the field by a boundary-value problem. It does not identify the joint $(u,m)$ system with an arbitrary initial-value semiflow. Classical regularity of $U$ is not asserted by this theorem.

\section{What compactness would suffice for noisy selection}
For each admissible repulsion family, let $U^\varepsilon$ solve exactly the uncorrected regularized master equation of the question. In shorthand its equation is
\begin{equation}\label{viscous}
 \partial_tU_i^\varepsilon+\varepsilon A:D^2U_i^\varepsilon
 +b^\phi(m,U^\varepsilon)\cdot DU_i^\varepsilon
 +2\varepsilon\sum_jm_j(\partial_iU_i^\varepsilon-\partial_jU_i^\varepsilon)
 -H_i(U^\varepsilon)+f_i+\sum_{j\ne i}\phi_\varepsilon(m_j)(U_j^\varepsilon-U_i^\varepsilon)=0,
\end{equation}
with $A_{jk}=m_j\delta_{jk}-m_jm_k$, terminal value $g$, and
$b_i^\phi=b_i+\phi_\varepsilon(m_i)-m_i\sum_j\phi_\varepsilon(m_j)$.

\begin{theorem}[Conditional interior selection]\label{compact}
Assume \eqref{short}, smooth $f,g$, and consider any collection of admissible repulsion families. Suppose for every family its solutions satisfy a common bound $|U_i^\varepsilon|\le M_0$, and for every compact $K\subset\operatorname{int}\Delta_d$ there is a bound, uniform in $\varepsilon$ and the families considered,
\[
 \sup_{[0,T]\times K}\bigl(|DU^\varepsilon|+|D^2U^\varepsilon|+|\partial_tU^\varepsilon|\bigr)\le C_K.
\]
Suppose $\phi_\varepsilon\to0$ locally uniformly on $(0,1]$ for each family. Then $U^\varepsilon$ converges locally uniformly on $[0,T]\times\operatorname{int}\Delta_d$ to the field of Theorem~\ref{contract}. In particular, its differences $Z_i^\varepsilon=U_i^\varepsilon-U_d^\varepsilon$ have a unique $L^1_{\mathrm{loc}}$ limit independent of the family, and the family-independence target in the question holds for these families.
\end{theorem}
\begin{proof}
On every compact interior cylinder, the bounds give equicontinuity of $U^\varepsilon$ in time and space. Arzel\`a--Ascoli and a diagonal extraction give locally uniform convergence to $W$. The spatial Hessian bound implies local equicontinuity of $DU^\varepsilon$ in space. To obtain time equicontinuity of these gradients, approximate each directional derivative by a difference quotient with step $h$ inside a slightly larger compact set: the Hessian bound gives error $O(h)$, while the time-Lipschitz bound on $U^\varepsilon$ gives variation $O(|t-s|/h)$. Choosing $h=\sqrt{|t-s|}$ proves local time equicontinuity. Thus, after extraction, the spatial gradients converge locally uniformly too.

In \eqref{viscous} the Hessian and extra-gradient terms multiplied by $\varepsilon$ vanish uniformly on interior compact sets, as do the repulsion terms. Integrate the equation between two times and pass to the limit. The resulting right-hand side is continuous, so $W$ is $C^1$ in time and satisfies the classical first-order master equation
\[
 \partial_tW_i+b(m,W)\cdot DW_i-H_i(W)+f_i=0,\qquad W(T,m)=g(m).
\]
It is locally $C^1$ in space and bounded by $M_0$.

Starting from any interior $(t_0,m_*)$, solve $\dot m=b(m,W(t,m))$. Local Lipschitz continuity gives local uniqueness. Boundedness of $W$ gives $m_i(t)\ge m_{*i}e^{-2M_0(d-1)T}$, so the trajectory remains in a compact interior region and continues to $T$. Along it, the chain rule shows $u_i(t)=W_i(t,m(t))$ solves \eqref{chars}. Theorem~\ref{contract} identifies $W(t_0,m_*)=U(t_0,m_*)$. Since every subsequential limit is identical, compactness yields convergence of the whole family. Any two families converge to this same limit; uniform convergence on compact cylinders implies the required $L^1$ convergence of differences on every compact measurable set.
\end{proof}

These derivative bounds are sufficient conditions, not consequences of the admissibility conditions in the question, which permit dependence on $\varepsilon$. This theorem handles a genuinely nonpotential class before singular gradients become an issue. It does not provide an intrinsic weak admissibility condition after such singularities.

\section{Why the potential gradient condition cannot be retained}
\begin{proposition}[A smooth nonpotential terminal field]\label{curl}
For $d=3$, choose $f=0$ and $g=(m_2,0,0)$. In the chart $x=m_1$, $y=m_2$, $m_3=1-x-y$, the terminal feedback-difference field is $Z=(y,0)$. It is not the gradient of any scalar distribution on any nonempty open subset of the simplex interior. Nevertheless it satisfies the one-sided Lipschitz estimate
\[
 (Z(z)-Z(z'))\cdot(z-z')\le\tfrac12|z-z'|^2.
\]
For sufficiently small $T$, Theorem~\ref{contract} applies to these data.
\end{proposition}
\begin{proof}
If $Z=DP$ distributionally, commutation of distributional mixed derivatives would imply $\partial_yZ_1=\partial_xZ_2$. Here these derivatives equal $1$ and $0$, respectively, a contradiction even on an arbitrarily small open set. For $z-z'=(a,b)$, the left side of the one-sided inequality is $ab\le(a^2+b^2)/2$. The data are smooth and bounded on the simplex, with $L_g\le1$ and $L_f=0$, so $8T<1$ suffices for \eqref{short}.
\end{proof}

Thus a one-sided Lipschitz bound alone is compatible with a nonzero curl. The scalar potential representation used for the known conservative weak theorem cannot simply be imposed in a class that must agree with nonpotential classical data. The remaining target is to find compactness and an admissibility/uniqueness principle that control both compression and this rotational part when derivatives are no longer uniformly bounded. No result here proves those estimates or regularization independence for the unrestricted admissible class. The full catalogue problem remains unresolved.

\begin{thebibliography}{9}
\bibitem{CD} P. Cardaliaguet and F. Delarue, \emph{Selected topics in mean field games}, Proc. ICM 2022, vol. 5, 3660--3703, EMS Press (2023), Theorem 4.6 and Section 4.4. \url{https://doi.org/10.4171/ICM2022/17}.
\bibitem{CeD} A. Cecchin and F. Delarue, \emph{Selection by vanishing common noise for potential finite state mean field games}, Comm. Partial Differential Equations 47 (2022), 89--168. Theorem 2.7 and Section 6 prove the potential conservative weak uniqueness result; they are not a nonpotential theorem. \url{https://arxiv.org/abs/2005.12153}.
\end{thebibliography}

\section{Completion Estimate}
20\%

\end{document}
